\documentclass[11pt]{article}

\usepackage[margin=1in]{geometry}
\usepackage{amsmath,amssymb,amsthm}
\usepackage[T1]{fontenc}
\usepackage{lmodern}
\usepackage[hidelinks]{hyperref}

\newtheorem{theorem}{Theorem}
\newtheorem{lemma}{Lemma}
\newtheorem{proposition}{Proposition}
\newtheorem{corollary}{Corollary}
\newtheorem{heuristic}{Heuristic}
\newtheorem{remark}{Remark}

\newcommand{\F}{\mathbb F}
\newcommand{\B}{\mathbb B}
\newcommand{\A}{\mathcal A}
\newcommand{\Deltaone}{\Delta_1}
\newcommand{\Deltazero}{\Delta_0}
\newcommand{\tauone}{\tau_1}
\newcommand{\tautwo}{\tau_2}
\newcommand{\tauthree}{\tau_3}
\newcommand{\Ctwo}{C_2}
\newcommand{\qline}{q_{\rm line}}
\newcommand{\qgen}{q_{\rm gen}}

\title{Experimental Theorems for the F1 / Slack-Two Resonance Window}
\author{RS-MCA experimental notes}
\date{June 18, 2026}

\begin{document}
\maketitle

\begin{abstract}
This note records restricted experimental theorems and heuristics extracted
from the Cycle 14--18 Fable-loop audits in the RS-MCA repository.  The main
new item is a divisibility gate for the non-coprime resonance branch in the
quadratic-extension, slack-two, three-co-support toy window.  These statements
are useful for Paper B's residue-line program, but they are not a proof of the
corrected MCA conjecture and have no direct protocol consequence.
\end{abstract}

\section{Status and Scope}

All results in this note live in the experimental folder and should be read as
restricted algebraic facts or proof targets.  We do not edit or supersede the
main papers A--D.

The common restricted setting is
\[
  B=\F_p,\qquad F=\F_{p^2},\qquad D=\F_p,\qquad
  t=\sigma=2,\qquad j=3,
\]
and we work off the tangent/global endpoint
\[
  R_0=\{\, [W]_E\wedge [B]_E=0\,\}.
\]
Here \(E\in F[X]\) has degree \(2\), \(b=[B]_E\), and
\[
  A=F[X]/E.
\]
The field ledgers remain distinct:
\[
  \qgen=p,\qquad \qline=p^2.
\]
The regime is sub-reserve, since \(\eta=\sigma/n=2/n\).  Thus none of the
following proves a corrected-reserve list, CA, MCA, line-decoding, SNARK, or
protocol theorem.

\section{Cycle 14 Normal Form}

The Cycle 14 audit reduces the \(t=2,j=3\) landing equation to two affine
forms in the elementary co-support parameters
\[
  \tau=(\tauone,\tautwo,\tauthree).
\]

\begin{lemma}[Affine wedge normal form]
\label{lem:affine-wedge}
Assume \(\{[W]_E,b\}\) is an \(F\)-basis of \(A\).  Then the landing equation
can be written in the form
\[
  \iota(\tau)=A_0(\tauone,\tautwo)-\tauthree [W]_E,\qquad
  \mu(\tau)=B_0(\tauone,\tautwo)-\tauthree b,
\]
where
\[
  A_0=p_1[W]_E+p_2 b,\qquad B_0=q_1[W]_E+q_2 b,
\]
and \(p_i,q_i\in F\) are affine-linear functions of
\((\tauone,\tautwo)\).  If the basis is normalized so
\([W]_E\wedge b=1\), then
\[
  \Delta=\iota\wedge\mu
    =(p_1-\tauthree)(q_2-\tauthree)-p_2q_1.
\]
\end{lemma}

\begin{proof}
The Cycle 14 quotient calculation gives the displayed affine forms.  Expanding
in the ordered basis \(\{[W]_E,b\}\), the coordinate vector of \(\iota\) is
\((p_1-\tauthree,p_2)\), while the coordinate vector of \(\mu\) is
\((q_1,q_2-\tauthree)\).  Their wedge determinant is therefore
\[
  (p_1-\tauthree)(q_2-\tauthree)-p_2q_1.
\]
\end{proof}

\section{Cycle 18 Component Split}

Choose a basis \(F=B+\alpha B\) with \(\alpha^2=\nu\in B\), and write
\[
  p_1=p_{10}+\alpha p_{11},\quad
  p_2=p_{20}+\alpha p_{21},\quad
  q_1=q_{10}+\alpha q_{11},\quad
  q_2=q_{20}+\alpha q_{21}.
\]
Split
\[
  \Delta=\Deltazero+\alpha\Deltaone,\qquad
  \Deltazero,\Deltaone\in B[\tauone,\tautwo,\tauthree].
\]

\begin{theorem}[Monic quadratic and linear alpha component]
\label{thm:component-split}
In the setting above,
\[
\begin{aligned}
\Deltazero={}&
  \tauthree^2-(p_{10}+q_{20})\tauthree
  +p_{10}q_{20}+\nu p_{11}q_{21}
  -p_{20}q_{10}-\nu p_{21}q_{11},\\
\Deltaone={}&
  -(p_{11}+q_{21})\tauthree
  +p_{10}q_{21}+p_{11}q_{20}-p_{20}q_{11}-p_{21}q_{10}.
\end{aligned}
\]
Consequently \(\Deltazero\) is monic quadratic in \(\tauthree\), while
\(\Deltaone\) has degree at most one in \(\tauthree\).
\end{theorem}

\begin{proof}
Substitute the coordinate expressions for \(p_i,q_i\) into the identity of
Lemma~\ref{lem:affine-wedge} and use \(\alpha^2=\nu\).  Collecting the
coefficient of \(1\) gives the displayed formula for \(\Deltazero\), and
collecting the coefficient of \(\alpha\) gives the displayed formula for
\(\Deltaone\).  The coefficient of \(\tauthree^2\) in \(\Deltazero\) is \(1\),
and the \(\alpha\)-component has no \(\tauthree^2\)-term.
\end{proof}

\begin{corollary}[Graph reduction outside the zero-alpha branch]
\label{cor:graph}
If \(\Deltaone\) is not the zero polynomial, write
\[
  \Deltaone=s(\tauone,\tautwo)\tauthree+h(\tauone,\tautwo).
\]
On the open locus \(s\ne0\), the common-zero condition
\(\Deltazero=\Deltaone=0\) is supported on the graph
\[
  \tauthree=-h/s.
\]
\end{corollary}

\begin{proof}
This is immediate from solving the linear equation \(\Deltaone=0\) for
\(\tauthree\) on \(s\ne0\).
\end{proof}

\section{New Divisibility-Gate Theorem}

The next theorem is the main new theorem extracted from the Cycle 18
reduction.  It is stronger than merely saying that the branch is a graph: it
shows that any two-dimensional family must satisfy an explicit divisibility
identity.

\begin{theorem}[Divisibility gate for the graph branch]
\label{thm:divisibility-gate}
Assume the restricted setting above and assume \(\Deltaone\ne0\).  Write
\[
  \Deltaone=s(\tauone,\tautwo)\tauthree+h(\tauone,\tautwo),
\]
and define the denominator-cleared graph polynomial
\[
  G(\tauone,\tautwo)
  =s(\tauone,\tautwo)^2
    \Deltazero\bigl(\tauone,\tautwo,-h(\tauone,\tautwo)/s(\tauone,\tautwo)\bigr).
\]
Then \(\deg G\le 4\).  If \(G\) is not the zero polynomial, then
\[
  \#\{\,\tau\in B^3:\Deltazero(\tau)=\Deltaone(\tau)=0\,\}=O(p).
\]
Consequently the corresponding number of distinct slopes satisfies
\[
  \Ctwo=O(p)=O(n).
\]
Therefore a growing-prime family with
\[
  \Ctwo=\Theta(p^2)=\Theta(\qline)
\]
in the nonzero-\(\Deltaone\) branch can occur only if \(G\equiv0\).
\end{theorem}

\begin{proof}
Because the \(p_i,q_i\) are affine-linear in \((\tauone,\tautwo)\), the
coefficient \(s\) has degree at most \(1\), and \(h\) has degree at most \(2\).
Since \(\Deltazero\) is monic quadratic in \(\tauthree\), clearing the
denominators after substituting \(\tauthree=-h/s\) gives a polynomial \(G\) of
degree at most \(4\).

On \(s\ne0\), the equation \(\Deltaone=0\) forces
\(\tauthree=-h/s\).  Hence every common zero on this open locus gives a base
pair \((\tauone,\tautwo)\) with \(G(\tauone,\tautwo)=0\).  If \(G\) is nonzero,
Schwartz--Zippel gives at most \(4p\) such base pairs, and each base pair gives
at most one value of \(\tauthree\).

It remains to count the exceptional locus \(s=0\).  There
\(\Deltaone=0\) also forces \(h=0\).  If \(s\) is a nonzero affine form, this
is contained in an affine line; if \(s\equiv0\), then \(h\) is a nonzero
polynomial of degree at most \(2\).  In either case there are \(O(p)\) base
pairs.  For each such pair, \(\Deltazero\) is a monic quadratic in
\(\tauthree\), so it gives at most two values of \(\tauthree\).  The total
common-zero count is \(O(p)\).

Finally, the number of distinct slopes is bounded by the number of landing
points \(\tau\).  Thus \(\Ctwo=O(p)\) whenever \(G\not\equiv0\), and any
\(\Theta(p^2)\) family in this branch must satisfy \(G\equiv0\).
\end{proof}

\begin{remark}
The identity \(G\equiv0\) says that, after clearing the \(s\)-denominator,
\(\Deltazero\) is divisible by the linear equation \(\Deltaone\) on the graph
branch.  This is now the first exact algebraic gate that a large slope family
must pass.
\end{remark}

\section{Nearby Proven Gates}

The same restricted lane contains two other useful gates.  They are included
here because they explain where the new theorem fits.

\begin{proposition}[Base-component complete-intersection gate]
\label{prop:complete-intersection}
Suppose \(\Deltazero,\Deltaone\in B[\tauone,\tautwo,\tauthree]\) have no common
surface component.  Then
\[
  \#V(\Deltazero,\Deltaone)(B)=O(p),
\]
and hence this branch contributes \(\Ctwo=O(p)\) slopes.
\end{proposition}

\begin{proof}
The two components have bounded degree.  If they share no surface component,
their common zero set in \(\mathbb A^3_B\) is a curve of bounded degree.  A
bounded-degree affine curve over \(\F_p\) has \(O(p)\) rational points, and
slopes are bounded by landing points.
\end{proof}

\begin{proposition}[Cycle 16 determinant gate]
\label{prop:det-gate}
Let \(Q(z_0,z_1)\) be the Cycle 15--16 determinant consistency polynomial for
the affine system
\[
  L_z(\tau)=\iota(\tau)-z\mu(\tau)=0,\qquad z=z_0+\alpha z_1.
\]
If \(Q\) is not identically zero, then in the restricted \(D=\F_p\),
\(t=\sigma=2\), \(j=3\), off-\(R_0\) window,
\[
  \Ctwo\le 4p=O(p)=O(n).
\]
\end{proposition}

\begin{proof}
Cycle 16 gives \(\deg Q\le4\) and shows that every landing slope must satisfy
\(Q(z_0,z_1)=0\).  If \(Q\ne0\), Schwartz--Zippel over \(B=\F_p\) gives at
most \(4p\) possible pairs \((z_0,z_1)\), hence at most \(4p\) slopes.
\end{proof}

\section{Slope Map and Remaining Heuristics}

On the graph branch, the landing equation \(\iota=z\mu\) gives
\[
  p_1-\tauthree=zq_1,\qquad p_2=z(q_2-\tauthree).
\]
Where \(q_1\ne0\), substituting \(\tauthree=-h/s\) yields
\[
  z(\tauone,\tautwo)=\frac{p_1+h/s}{q_1}.
\]
Equivalently, eliminating \(\tauthree\) recovers the Cycle 14 quadratic
\[
  q_1z^2-(p_1-q_2)z-p_2=0,
\]
so the residual map is controlled by
\[
  (\tauone,\tautwo)\longmapsto [\,q_1:(p_1-q_2):p_2\,]\in\mathbb P^2(F).
\]

\begin{heuristic}[Exceptional identity should be rare]
For generic source data, the polynomial \(G\) should not vanish identically.
Thus the graph branch should usually be curve-sized, and the first possible
large-slope families should come from algebraically special data satisfying
the divisibility identity \(G\equiv0\) or from the zero-alpha branch
\(\Deltaone\equiv0\).
\end{heuristic}

\begin{heuristic}[What remains after the gate]
If \(G\equiv0\), the problem is no longer to show that the common-zero set is
small.  It may be surface-sized.  The remaining question is whether the
projective coefficient map has one-dimensional image on every source-valid
surface, giving \(\Ctwo=O(p)\), or whether some growing-prime source-valid
family has two-dimensional image and \(\Ctwo=\Theta(p^2)\).
\end{heuristic}

\section{Next Experimental Checks}

A useful scanner should compute, for each source-valid candidate:
\[
  \Deltaone=s\tauthree+h,\qquad
  G=s^2\Deltazero(\tauone,\tautwo,-h/s),
\]
then record whether \(G\equiv0\), whether \(\Deltaone\equiv0\), the
exceptional locus \(s=h=0\), the image size of
\[
  [\,q_1:(p_1-q_2):p_2\,],
\]
and the direct slope count over distinct split triples \(T\subset\F_p\).

The certificate should keep at least:
\[
\begin{gathered}
  p,\ \qgen,\ \qline,\ E,\ B,\ \text{seed},\ \text{stratum},\
  R_0\text{-status},\\
  \Deltaone\equiv0,\quad G\equiv0,\quad
  \#\{s=h=0\},\quad \#\operatorname{image},\quad
  \Ctwo,\quad \#\{T\text{ split distinct}\}.
\end{gathered}
\]

\section{Non-Claims}

This note does not claim:
\begin{itemize}
\item a proof of the corrected MCA conjecture;
\item a positive theorem above corrected reserve;
\item a \(\qgen\) local-limit theorem;
\item a list, CA, MCA, line-decoding, SNARK, or protocol consequence;
\item a large \(\Theta(p^2)\) counterpacket.
\end{itemize}

The mathematical value is narrower: the possible large branch in this
restricted window must satisfy an explicit algebraic identity before it can
even be surface-sized.

\end{document}
